% !TeX program = pdflatex
https://healthacademy2026.com/\documentclass[10pt]{beamer}

\usetheme{Madrid}
\setbeamertemplate{navigation symbols}{}
\setbeamertemplate{footline}[frame number]

\usepackage[spanish,es-nodecimaldot]{babel}
\usepackage{amsmath,amssymb}

\title[Financiamiento en Salud]{Financiamiento de los Sistemas de Salud}
\subtitle{Riesgo, contrato social y economía generacional}
\author{Curso de Economía Poblacional}
\institute{Tecnológico de Monterrey}
\date{\today}

\begin{document}

%------------------------------------------------
\begin{frame}
\titlepage
\end{frame}

%------------------------------------------------
\begin{frame}{Ruta Parte I}
\begin{enumerate}
\item Salud como riesgo y bienestar (Shafik)
\item Seguro y pooling
\item Beveridge vs Bismarck
\item Informalidad y decisión microfundada
\item Externalidades y bien público
\item Transición demográfica y epidemiológica
\item Economía generacional y NTA
\end{enumerate}
\end{frame}

%================================================
\section{Salud como Riesgo}

\begin{frame}{Salud como riesgo económico}
Ingreso individual: $y$

Shock de salud:
\[
m \in \{0, M\}, \quad \Pr(M)=p
\]

Utilidad esperada sin seguro:
\[
EU = (1-p)u(y) + p\,u(y-M)
\]

\begin{itemize}
\item La enfermedad es un shock incierto.
\item Con aversión al riesgo, la variabilidad reduce bienestar.
\end{itemize}
\end{frame}

%------------------------------------------------
\begin{frame}{Shafik: riesgo y contrato social}
\begin{itemize}
\item La enfermedad no es elegida.
\item El mercado puro genera exclusión o empobrecimiento.
\item El financiamiento colectivo compra estabilidad.
\end{itemize}

\bigskip
\centering
\emph{El sistema de salud es una institución para compartir riesgo biológico.}
\end{frame}

%------------------------------------------------
\begin{frame}{Seguro y bienestar}
Prima actuarialmente justa:
\[
\pi = pM
\]

Consumo asegurado:
\[
c = y - \pi
\]

Si $u''<0$:
\[
u(y-\pi) > EU_{\text{sin seguro}}
\]

\bigskip
El pooling aumenta bienestar cuando hay aversión al riesgo.
\end{frame}

%================================================
\section{Modelos Institucionales}

\begin{frame}{Tres funciones del financiamiento}
\begin{enumerate}
\item Recaudación
\item Pooling
\item Compra estratégica
\end{enumerate}

\bigskip
La arquitectura institucional define cómo se comparten riesgos.
\end{frame}

%------------------------------------------------
\begin{frame}{Modelo Beveridge (1942)}
Financiamiento vía impuestos generales:

\[
t \sum_i y_i = \sum_i p_i M_i
\]

\begin{itemize}
\item Universalidad basada en ciudadanía.
\item Pool único.
\item Salud como bien social.
\end{itemize}
\end{frame}

%------------------------------------------------
\begin{frame}{Modelo Bismarck (1883)}
Contribución laboral:

\[
\sum_{i \in F} \tau w_i = \sum_{i \in F} p_i M_i
\]

\begin{itemize}
\item Cobertura ligada al empleo formal.
\item Multiplicidad de fondos.
\item Solidaridad dentro del grupo asegurado.
\end{itemize}
\end{frame}

%------------------------------------------------
\begin{frame}{Comparación conceptual}
\begin{itemize}
\item Beveridge: universalidad vía impuestos.
\item Bismarck: universalidad vía formalidad laboral.
\item En economías con alta informalidad, Bismarck puro no es viable.
\end{itemize}
\end{frame}
%------------------------------------------------
\begin{frame}{Decisión formal/informal: parámetros del modelo}

Shock de salud:
\[
m \in \{0,M\}, \quad \Pr(M)=p
\]

\begin{itemize}
\item $\tau$: tasa de contribución laboral.
\item $\kappa$: costo fijo de formalidad (regulación, cumplimiento).
\item $\alpha$: fracción del gasto cubierta en formalidad.
\item $\beta$: fracción cubierta en informalidad.
\item $M$: costo del tratamiento.
\item $p$: probabilidad del shock.
\end{itemize}

\bigskip
\emph{El individuo compara el costo presente de contribuir contra la protección futura.}

\end{frame}

%-----------------------------------------

%================================================
\section{Informalidad y Microfundamentos}

\begin{frame}{Decisión formal/informal}
Formal:
\[
EU_F = (1-p)u(y(1-\tau)-\kappa)
+ p\,u(y(1-\tau)-\kappa-(1-\alpha)M)
\]

Informal:
\[
EU_I = (1-p)u(y)
+ p\,u(y-\beta M)
\]

Elige formal si:
\[
EU_F \ge EU_I
\]
\end{frame}

%------------------------------------------------
\begin{frame}{Informalidad estructural en AL}
\begin{itemize}
\item Alta proporción fuera del sistema contributivo.
\item Contribuciones elevadas reducen incentivos a formalizar.
\item Resultado: pools pequeños y fragmentación.
\end{itemize}

\bigskip
\emph{En América Latina, Bismarck nunca es puro.}
\end{frame}

%================================================
\section{Externalidades}

\begin{frame}{Prevención y externalidades}
\[
p_i = p(e_i,\bar e)
\]

\[
\frac{\partial p}{\partial \bar e}<0
\]

\begin{itemize}
\item Las decisiones individuales afectan a terceros.
\item Vacunación y vigilancia son bienes públicos.
\item El mercado subprovee prevención.
\end{itemize}
\end{frame}

%================================================
\section{Transición Demográfica y Epidemiológica}

\begin{frame}{Estructura del gasto}
\[
G = \sum_x N(x)E[c(x)]
\]

\[
E[c(x)] = p_C(x)M_C(x) + p_N(x)M_N(x)
\]

\begin{itemize}
\item Transición epidemiológica: menos contagiosas, más crónicas.
\item Las crónicas implican tratamiento prolongado.
\end{itemize}
\end{frame}

%------------------------------------------------
\begin{frame}{Demografía y presión estructural}
\[
\frac{\partial G}{\partial N(65+)} > 0
\]

\begin{itemize}
\item Envejecimiento aumenta la demanda.
\item Cambia la naturaleza del gasto.
\end{itemize}

\bigskip
La demografía redefine el financiamiento óptimo.
\end{frame}

%================================================
\section{Economía Generacional}

\begin{frame}{Salud y NTA}
\[
LCD(x)=C(x)-Y^L(x)
\]

\[
LCD(x)=\tau_H(x)+\text{otros flujos}
\]

\begin{itemize}
\item La salud pública es transferencia en especie.
\item Concentrada en edades avanzadas.
\end{itemize}
\end{frame}

%------------------------------------------------
\begin{frame}{Contrato intergeneracional}
\[
\sum_x N(x)\tau_y(x)
=
\sum_x N(x)\tau_H(x)
\]

\begin{itemize}
\item Trabajadores financian.
\item Mayores reciben.
\item La estructura etaria altera ambos lados.
\end{itemize}
\end{frame}

%------------------------------------------------
\begin{frame}{Cierre Parte I}
\begin{itemize}
\item Salud = riesgo compartido.
\item Instituciones determinan pooling.
\item Informalidad limita cobertura.
\item Transición epidemiológica transforma el gasto.
\item El financiamiento es también contrato entre generaciones.
\end{itemize}


\end{frame}

\end{document}
